\chapter{2025年上海卷(春)}

\section{填空题}

\begin{problem}
已知集合 \(A=\{x\mid x>0\}\),\(B=\{-1,0,1,2\}\),则 \(A\cap B=\fillinblank{}\).
\end{problem}

\begin{problem}
不等式 \(\dfrac{x}{x-1}<0\) 的解集为\fillinblank{}.
\end{problem}

\begin{problem}
已知复数 \(z=\dfrac{2+\i}{\i}\),其中 \(\i\) 为虚数单位,则 \(|z|=\fillinblank{}\).
\end{problem}

\begin{problem}
已知 \(\bs{a}=(2,1)\),\(\bs{b}=(1,x)\),若 \(\bs{a}\parallel\bs{b}\),则 \(x=\fillinblank{}\).
\end{problem}

\begin{problem}
已知 \(\tan\alpha=1\),则 \(\cos\left(\alpha+\dfrac{\pi}{4}\right)=\fillinblank{}\).
\end{problem}

\begin{problem}
已知 \(\left(x+\dfrac{m}{x}\right)^6\) 的展开式中常数项为 \(20\),则实数 \(m\) 的值为\fillinblank{}.
\end{problem}

\begin{problem}
已知 \(\{a_n\}\) 是首项为 \(1\),公差为 \(1\) 的等差数列,\(\{b_n\}\) 是首项为 \(1\),公比为 \(q(q>0)\) 的等比数列. 若数列 \(\{a_nb_n\}\) 的前三项和为 \(2\),则 \(q=\fillinblank{}\).
\end{problem}

\begin{problem}
关于 \(x\) 的方程 \(|x-1|+|\pi-x|=\pi-1\) 的解集为\fillinblank{}.
\end{problem}

\begin{problem}
已知 \(P\) 是一个圆锥的顶点,\(PA\) 是母线,\(PA=2\),该圆锥的底面半径是 \(1\). \(B\),\(C\) 分别在圆锥的底面上,则异面直线 \(PA\) 与 \(BC\) 所成角的最小值为\fillinblank{}.
\end{problem}

\begin{problem}
已知双曲线 \(\dfrac{x^2}{a^2}-\dfrac{y^2}{6-a^2}=1(a>0)\) 的左,右焦点分别为 \(F_1\),\(F_2\). 通过 \(F_2\) 且倾斜角为 \(\dfrac\pi3\) 的直线与双曲线交于第一象限的点 \(A\),延长 \(AF_2\) 至 \(B\) 使得 \(AB=AF_1\). 若 \(\triangle BF_1F_2\) 的面积为 \(3\sqrt6\),则 \(a=\fillinblank{}\).
\end{problem}

\begin{problem}
如图,正方形 \(ABCD\) 是一块边长为 \(4\) 的工程用料,阴影部分是被腐蚀的区域,其余部分完好,曲线 \(MN\) 为以 \(AD\) 为对称轴的抛物线的一部分,且 \(DM=DN=3\). 工人师傅现要从完好的部分中截取一块矩形原料 \(BQPR\),当其面积有最大值时,\(AQ\) 的长为\fillinblank{}.

\bitmapfigure[max width=0.42\linewidth]{img/2025/shanghai_spring/q11_fig1.png}
\end{problem}

\begin{problem}
在平面中,\(\bs{e}_1\) 和 \(\bs{e}_2\) 是互相垂直的单位向量,向量 \(\bs{a}\) 满足 \(|\bs{a}-4\bs{e}_1|=2\),向量 \(\bs{b}\) 满足 \(|\bs{b}-6\bs{e}_2|=1\),求 \(\bs{b}\) 在 \(\bs{a}\) 方向上的数量投影的最大值\fillinblank{}.
\end{problem}

\section{单选题}

\begin{problem}
如图,\(ABCD-A_1B_1C_1D_1\) 是正四棱台,则下列各组直线中属于异面直线的是(\quad).

\bitmapfigure[max width=0.48\linewidth]{img/2025/shanghai_spring/q13_fig1.png}
\choices
    {\(AB\) 和 \(C_1D_1\)}
    {\(AA_1\) 和 \(CC_1\)}
    {\(BD_1\) 和 \(B_1D\)}
    {\(A_1D_1\) 和 \(AB\)}
\end{problem}

\begin{problem}
幂函数 \(y=x^\alpha\) 在 \((0,+\infty)\) 上是严格减函数,且经过 \((-1,-1)\),则 \(\alpha\) 的值可能是(\quad).
\choices
    {\(-\dfrac23\)}
    {\(-\dfrac13\)}
    {\(\dfrac13\)}
    {\(3\)}
\end{problem}

\begin{problem}
有一四边形 \(ABCD\),对其四边 \(AB\),\(BC\),\(CD\),\(DA\) 按顺序分别抛掷一枚质量均匀的硬币:如硬币正面朝上,则将其擦去;如硬币反面朝上,则不擦去. 最后,以 \(A\) 为起点沿着尚未擦去的边出发,可以到达 \(C\) 点的概率为(\quad).
\choices
    {\(\dfrac12\)}
    {\(\dfrac7{16}\)}
    {\(\dfrac14\)}
    {\(\dfrac3{16}\)}
\end{problem}

\begin{problem}
已知 \(a\in\R\),不等式 \(\left[\tan\left(\dfrac\pi6x\right)-a\right]\left[\tan\left(\dfrac\pi6x\right)-a-1\right]<0\) 在 \((0,2025)\) 中的整数解有 \(m\) 个. 关于 \(m\) 的个数,以下不可能的是(\quad).
\choices
    {\(0\)}
    {\(338\)}
    {\(674\)}
    {\(1012\)}
\end{problem}

\section{解答题}

\begin{problem}
在三棱锥 \(P-ABC\) 中,平面 \(PAC\perp\) 平面 \(ABC\),\(PA=AC=CP=2\),\(AB=BC=\sqrt2\).

\begin{enumerate}
\item 若 \(O\) 是棱 \(AC\) 的中点,证明:\(BO\perp\) 平面 \(PAC\),并求三棱锥 \(B-OPA\) 的体积;
\item 求二面角 \(B-PC-A\) 的大小.
\end{enumerate}

\FigureLayoutDeclare{img/2025/shanghai_spring/q17_fig1.png}{4.0cm}{0bp 0bp 0bp 0bp}
\bitmapfigure[max width=0.42\linewidth]{img/2025/shanghai_spring/q17_fig1.png}
\end{problem}

\begin{problem}
在 \(\triangle ABC\) 中,角 \(A\),\(B\),\(C\) 所对的边分别为 \(a\),\(b\),\(c\),且 \(c=5\).
\begin{enumerate}
\item 若 \(\dfrac{a}{4b}=\dfrac{\sin B}{\sin A}\),\(C=\dfrac\pi2\),求 \(a\);
\item 若 \(ab=20\),求 \(\triangle ABC\) 面积的最大值.
\end{enumerate}
\end{problem}

\begin{problem}
甲,乙是两个体育社团的小组,给出两组组员身高的茎叶图(单位:厘米),以身高的百位数和十位数作为``茎''排列在中间,个位数作为``叶''分列在两边.

\begin{center}
\begin{tabular}{c|c|c}
甲 & 茎 & 乙\\
\hline
 & \(15\) & \(9\)\\
\(7\),\(7\),\(5\),\(5\),\(4\) & \(16\) & \(0\),\(3\),\(5\),\(5\),\(6\),\(7\),\(8\),\(8\)\\
\(8\),\(5\),\(4\),\(3\),\(2\),\(2\) & \(17\) & \(2\)\\
\(3\) & \(18\) &
\end{tabular}
\end{center}

\begin{enumerate}
\item 分别求甲,乙两组组员身高的第 \(60\) 百分位数;
\item 从甲,乙两组各选取一个组员,求两人身高均在 \(170\text{ 厘米}\) 以上的概率;
\item 为使两组人数相同,从甲组中调派一个队员到乙组. 是否存在甲组的一个组员,将他调派至乙组后,甲,乙两组的平均身高都增大?
\end{enumerate}
\end{problem}

\begin{problem}
在平面直角坐标系中,已知曲线 \(\Gamma:\dfrac{x^2}{4}+y^2=1(y\ge0)\),点 \(P\),\(Q\) 分别为 \(\Gamma\) 上不同的两点,\(T(t,0)\).
\begin{enumerate}
\item 求 \(\Gamma\) 所在椭圆的离心率;
\item 若 \(T(1,0)\),\(Q\) 在 \(y\) 轴上,且 \(T\) 到直线 \(PQ\) 的距离为 \(\dfrac{\sqrt5}{5}\),求 \(P\) 的坐标;
\item 是否存在 \(t\),使得 \(\triangle TPQ\) 是以 \(T\) 为直角顶点的等腰直角三角形?若存在,求 \(t\) 的取值范围;若不存在,请说明理由.
\end{enumerate}

\bitmapfigure[max width=0.42\linewidth]{img/2025/shanghai_spring/q20_fig1.png}
\end{problem}

\begin{problem}
已知函数 \(y=f(x)\) 的定义域是 \(D\). 对于 \(t\in D\),定义集合 \(S_f(t)=\{x\mid f(x)\ge f(t)\}\).
\begin{enumerate}
\item \(f(x)=\log_2x\),求 \(S_f(16)\);
\item 对于集合 \(A\),若对任意 \(x\in A\) 都有 \(-x\in A\),则称 \(A\) 是对称集. 若 \(D\) 是对称集,证明:``函数 \(y=f(x)\) 是偶函数''的充要条件是``对任意 \(t\in D\),\(S_f(t)\) 是对称集'';
\item 若 \(x\in\R\),\(f(x)=\e^x-\dfrac12mx^2\). 求 \(m\) 的取值范围,使得对于任意 \(t_1<t_2\in D\),都有 \(S_f(t_2)\subseteq S_f(t_1)\).
\end{enumerate}
\end{problem}

